Ten structures were selected from the database for detailed fracture analysis, chosen for what they teach rather than for their strength alone: the two reference states (pristine zigzag graphene and the 20\,\AA{} precracked sheet), the strongest porous architecture, the most delocalised and the most localised failures, the key hierarchical design, the flaw-halo design, the weakest member of the matched-porosity hierarchy ladder, and the two holdout designs that changed the conclusions (the aligned two-level mesh and the 20-degree slit array). Two further categories, the largest work to failure (the 0.20-porosity parallel slit array, 3.25\,J/m$^2$, whose mechanism is that of the strongest design and which appears in Fig.~\ref{fig:comparison_slit_orientation_echelon}) and the load-parallel-vein mesh of Stage 4 (Fig.~\ref{fig:discriminating_experiments}), were analysed in the same way; their panels and movies are included in the archive but omitted here for length.

For each structure the progression panel shows event-selected frames rather than equally spaced ones: the relaxed state, the last frame before damage, the first bond loss, the peak load, the largest single propagation event, the first post-peak state and the final frame, each labelled with strain, 2D stress and the cumulative number of broken bonds. Atoms are coloured by their energy change relative to the relaxed state (the per-atom energies of the potential), and atoms that have lost a bond are circled. The before/after panels show the relaxed, peak-load and final configurations side by side. Movies of the same runs, with a synchronised stress--strain inset, are in \texttt{movies/} and in the ``Fracture progression'' panel of the app. The interpretation below each panel is a statement about the model; the ``interpretation'' of a structure as a mechanism was written after inspecting the frames and cross-checked against the matched comparisons of Stages 4 and 5.
